\documentclass[11pt]{letter}
\usepackage[margin=2.5cm]{geometry}

\signature{David Alarc\'on\\Universidad Pablo de Olavide, Sevilla, Spain\\dalarub@upo.es}
\address{David Alarc\'on\\Universidad Pablo de Olavide\\Ctra.\ de Utrera, km 1\\41013 Sevilla, Spain}

\begin{document}
\begin{letter}{Editorial Board\\Journal of Number Theory}

\opening{Dear Editors,}

I am submitting the manuscript ``Ab initio derivation of the Berry-Keating
correction coefficient from the Conrey-Snaith triple correlation formula''
for consideration in the \textit{Journal of Number Theory}.

The gap ratio statistic of Riemann zeta zeros converges to GUE as
$\langle r\rangle(T) = R_\infty + c/\log^2 T$ with $c = 1.245 \pm 0.040$
(measured empirically in a companion paper). Here we derive $c$ from first
principles using the Conrey-Snaith formula for the triple correlation
function $R_3(v_1, v_2; L)$, obtaining $c_{\mathrm{teo}} = 1.23$ (99\%
of the empirical value). This is the first ab initio calculation of a
Berry-Keating correction coefficient for any spacing statistic.

The paper also documents 29 failed perturbative approaches, identifying
the structural obstruction ($\mathrm{sinc}(1) = 0$) that prevents any
kernel-based derivation.

This work is part of a companion series: Paper~A (\textit{Comm.\ Math.\
Phys.}) presents the data and mechanism, and Paper~C (\textit{Annals of
Mathematics}) proves that the convergence implies the Riemann Hypothesis.

The manuscript has not been published elsewhere and is not under
consideration by another journal.

\closing{Sincerely,}

\end{letter}
\end{document}
